% =====================================================================
%  Harald Høffding, Review of Viktor Rydberg, Ny Testamentes Lærdomme om Forudtilværelsen og de sidste Ting
%  Translated from the Danish: Ude og Hjemme, vol. 5 (1881--82), no. 224 (15 January 1882), pp. 196--197.
%
%  TRANSLATION -- GENERATED; DO NOT EDIT BY HAND.  Built by ebook-method/rydberg-anmeldelse-1882/tr_build.py --write from
%  tr_template.tex and tr/RY.texfrag (translated 2026-10-17 in the main session, from the transcription read whole at the
%  image: texts/hoeffding/rydberg-anmeldelse-1882/transcription.tex); [N] in the margin = the page of the original.
%  Terms as the earlier Høffding translations (ebook-method/*/GLOSSARY.md): Forestilling = idea (religious ideas),
%  Forestillingsformer = forms of representation, Forestillingskreds = circle of ideas; Forudtilværelse = pre-existence;
%  Udvikling = development; Erfaringsvidenskab = empirical science; Aand = spirit.  Rydberg's titles stay Swedish.
%  The two misprints of the print are rendered in their evident sense (a `% print:' line above each).
% =====================================================================
\documentclass[12pt,a4paper,openany]{book}
\usepackage[utf8]{inputenc}
\usepackage[T1]{fontenc}
\usepackage{libertinus}
\usepackage{libertinust1math}
\usepackage[german,french,english]{babel}
\usepackage{textalpha}
\usepackage[protrusion=true,expansion=true]{microtype}
\usepackage{setspace}
\usepackage[margin=1.3in]{geometry}
\usepackage{fancyhdr}
\usepackage{hyperref}
\hypersetup{pdftitle={Review of Viktor Rydberg, Ny Testamentes Lærdomme}, pdfauthor={Harald Høffding}, hidelinks}
\pagestyle{fancy}
\fancyhf{}
\fancyhead[LE,RO]{\thepage}
\fancyhead[C]{\textit{H. Høffding: Review of Viktor Rydberg, Ny Testamentes Lærdomme (1882)}}
\fancypagestyle{plain}{\fancyhf{}\fancyhead[LE,RO]{\thepage}}
\setstretch{1.3}
\usepackage{marginnote}
\usepackage{graphicx}
\DeclareUnicodeCharacter{0254}{\reflectbox{c}}   % ɔ (id est)

\newcommand{\opage}[1]{\marginnote{\footnotesize\textit{[#1]}}}
\newcommand{\signature}[1]{\par\smallskip\hfill #1\hspace*{2em}\par}
\pdfstringdefDisableCommands{\renewcommand{\footnote}[2][]{}\renewcommand{\opage}[1]{}}

\renewcommand{\thefootnote}{*}   % the print marks its one note *)
\makeatletter\renewcommand{\@makefnmark}{\hbox{\textsuperscript{\@thefnmark)}}}\makeatother
\begin{document}

\begin{titlepage}
\begin{center}
{\Huge \emph{Ude og Hjemme}}\\[0.4em]{\large [Abroad and at Home]}\\[0.6em]{\large A Nordic Illustrated Weekly}\\[0.6em]{\small edited by Otto Borchsenius and F. Hendriksen}\\[0.6em]{\small Fifth Volume, 1881--1882}\\[0.4em]{\small Copenhagen: F. Hendriksen}\\[2em]
{\small Translated from the Danish}
\end{center}
\end{titlepage}
\clearpage

% --- p. 196 ---
\setcounter{footnote}{0}%
\opage{196}\begin{center}{\sffamily\large Viktor Rydberg: Ny Testamentes Lærdomme om Forudtilværelsen og de sidste Ting.}\\[0.2em]{\small [The New Testament's Teachings on Pre-existence and the Last Things]}\\[0.3em]{\small Danish translation by Cand.\ theol.\ \textbf{\sffamily J. P. V. Petersen.}  Copenhagen 1881.  (Andr. Schous Forlag.)}\end{center}
% (p. 196, column 1)

While Viktor Rydberg is well known in wide circles as a poet, translator and art historian, fewer here at home know that he also occupies an eminent place as a historian of religion and is one of the very few scholars in the Nordic countries who, without belonging to the ranks of ecclesiastical orthodoxy, have devoted themselves with zeal and enthusiasm to the study of the biblical writings. Instead of consciously or unconsciously taking the prevailing doctrine of the Church as the basis for interpreting the books of the New Testament, Rydberg demands that one see them in the light of the time in which they came into being, and let the thoughts they contain unfold in full freedom, disregarding the way in which later ages have arranged them to suit themselves.
% print: Fortolkning,. (misprint: a comma and a stop)
It is, in other words, a purely historical interpretation that Rydberg advocates. But however much he keeps his own opinions and sympathies out of his research, a merely historical understanding is nevertheless not his final aim. It is his conviction that the deeper one penetrates to the wellspring of the great religious movement from which those ancient writings arose, the more one will also find means of freeing religious feeling from the strife and doubt that have spread through the spiritual world as a consequence of the collision between ecclesiastical theology and scientific thought. The picture of Christ that the New Testament writings give us, when they are understood in their own light, contains in his conviction none of the contradictions with which the doctrine of the Church is burdened, but is able to satisfy both thought and the heart. By pushing aside ecclesiastical orthodoxy in favor of the original New Testament teaching, which also prevailed in the Church before Christianity became the religion of the state, the author wishes to contribute his mite to the work of continuing the Reformation and building ``a Church that encourages reason to inquire, just as it encourages the heart to love, since it seeks to realize both truth and love, and through both, freedom''.
% (p. 196, column 2)
Preliminary work toward such a continuation of what was begun with Luther's Reformation, but halted all too early, the author finds here and there throughout the Christian world. In his own fatherland ``the pious and dauntless Ignell and the great thinker Boström'' in particular have worked in this direction.

The present work, which is connected with the author's larger work ``Bibelns lära om Kristus'', which appeared in a fourth printing in 1880, seeks to demonstrate that the belief in man's pre-existence---which the ancient world felt no less need to assume than an existence after death, and which in particular pervaded the whole religious circle of ideas of the time in which the New Testament writings came into being---that this belief is also the teaching of the New Testament. Here the author relies especially on the Gospel of John, the Pauline epistles and the Epistle to the Hebrews. The human personalities that exist in the world of sense in separation and imperfection existed in original unity in the eternal ideal, which the Gospel of John calls the Logos and Paul the heavenly Adam. According to the doctrine of the Messiah, it is the task of the historical Christ to unite what has been separated and to gather what has been scattered. The world of sense is a school through which the original, ideal unity is to be realized again. According to the New Testament, historical development ends in a renewal of the world, in which each single personality takes its place in the ideal world of spirits. The doctrine of an eternal damnation, which has prevailed in the Church since the fifth century, the author declares unbiblical. He does not deny that ``the doctrine of retribution, on which the Mosaic Law rests, and which even the greatest of Israel's prophets could not wholly overcome, also casts its shadow into the New Testament''; but he finds it at odds with the prevailing line of thought in the New Testament writings. Especially
% --- p. 197 ---
\setcounter{footnote}{0}%
% (p. 197, column 1)
\opage{197}interesting is his demonstration of the misunderstandings that have arisen from incorrect translation and interpretation of biblical expressions (as when ``Hades'' is simply translated by ``Hell'', ``aionios'' by ``eternal''). --- I shall not enter here upon a closer discussion of these questions. Such a discussion would presuppose a special exegetical learning that I do not possess. I can only say this much, that the author's presentation and interpretation seem unforced and lead the reader into a circle of ideas that seems to accord with the time and the circumstances under which the New Testament writings came into being. The author's historical sense and great love of truth\footnote[1]{A theological opponent in Sweden has not been able to go further than to grant that Rydberg ``seems (!) to seek the truth with a certain (!) sincerity''.} forbid him to rest before he has reached the primitive layers of the soil he is examining.

One cannot read an investigation like the present one without thinking of what will happen if the education of future clergymen should become still less scientific and more like that of a seminary than it is at present. If Rydberg is right, false translation of biblical expressions has led to unbiblical doctrines in the Church. But the theologian without linguistic and critical training must keep to the translation and is excluded from correcting such errors, indeed even from judging a contribution such as Rydberg's. The decision whether Rydberg is right does not depend on whether he is a ``believer'' or not, but on whether he is right in his linguistic and historical interpretation. ---

Viktor Rydberg's own religious standpoint coincides, as one sees from several statements in his writings, with what he regards as genuinely biblical only in so far as the biblical ideas of pre-existence, of temporal and sensuous existence and of the renewal of the world can serve as symbols for religious ideas that in themselves point beyond the forms of time and space. Rydberg adheres, that is, to the idealistic philosophy prevailing in Sweden, according to whose doctrine the divine cannot be subject to becoming and development. Christianity's conception of the life of the world as a great historical process of development, in which the Deity itself works and suffers along, can have only symbolic value for such a standpoint. Even Hegel, after all, is too empirical for the Swedish idealists. Here it seems to me that one might have expected a freer view from the Swedish Lessing. I do not reproach him for seeing a poetic symbolism in the biblical ideas and for being unable to accept the literal conception of them. That is how I myself stand toward these venerable traditions, finding, moreover, the sum of biblical wisdom not in any of the dogmatic ideas whatever, but in the proclamation of the gospel of unconditional love of humanity and in the beatitude of those who hunger and thirst after righteousness. But I believe that Rydberg could have gone a step further and acknowledged it as significant that the Christian foundational writings, in their own way, acknowledge development as a law of the world. When modern science arrives at the same thought from so many sides, this would surely be a fruitful starting point for anyone who wishes to seek out what is common behind the conflicting oppositions.
% print: Udganspunkt (misprint for Udgangspunkt)
An idealism that regards time and development as mere illusions shuts itself out from all fruitful intercourse with the stirrings of real life. It is possible that concepts such as time and development, like all other concepts, have validity only from a human standpoint; but we simply cannot place ourselves at any other. Critical insight into the connection of our fundamental concepts with our organization and our place in the world is not the same as a conviction that they are illusory; on the contrary: if our fundamental concepts really are connected with our nature, then they also have all the validity that we can attribute to concepts at all. --- If Rydberg did not start from so abstract an idealism, he would also follow the more recent endeavors in the psychological and philosophical field with greater sympathy. There is evidently a gap here in his rich and many-sided cultivation. While
% (p. 197, column 2)
his thought is honed for the understanding and treatment of poetic, religious and speculative forms of representation, he has hardly made himself at home, with the same diligence and sympathy, in the method of the empirical sciences and its significance for philosophy as well. How foreign he is to the purely empirical (inductive) way of looking at things one gets an impression of when he occasionally mentions the younger Fichte and Carus as researchers who ``independently of metaphysical theories have attempted to penetrate the problem of anthropology by the path of natural science''. The men named do, to be sure, very often appeal to ``experience''; but they then count as such all manner of unclear mysticism and spiritualist fantasy, when it suits their opinions. The true empirical method was far from them. (Cf.\ Lotze's polemic against the younger Fichte.)

These questions are touched upon only by way of hint in the book, although the hints make the philosophical reader curious. It would surely be of great significance for those in the North who, each in his own way, embrace a freer religious conviction, if a man such as Viktor Rydberg would give us an independent treatment of the religious problem. With his clear thought, his splendid language, and above all his warm feeling for all that is noble and true, he has important qualifications for solving such a task. His untiring research would then probably also lead him to find a good deal of usable ore even in those of the philosophical endeavors of the present day in which, as it seems, he has so far seen only slag.
\signature{\textsf{\bfseries Harald Høffding.}}

\end{document}
